\section{Related Work}
\subsection{Graph Retrieval-Augmented Generation}
RAG augments LLMs with external evidence to improve factuality of knowledge-intensive tasks~\citep{rag}. Building on this, GraphRAG is an advance paradigm that explicitly models structural relations among entities, thereby capturing relational semantics, contextual dependencies and structural knowledge integration~\citep{graphragsurvey,graphrag}.
Early work~\citep{raptor, graphrag} emphasize hierarchical summarization and global information integration, but they insufficiently leveraged the fine-grained structural information. LightRAG~\citep{lightrag} advanced this direction by incorporating graph structures into both indexing and retrieval.
Recent efforts in graph KB construction introduce diverse structural representations, such as the design of heterogeneous and lightweight graph structures~\citep{minirag, noderag}, the extension to hypergraphs that capturing higher-order relational dependencies~\citep{hypergraphrag, hyperrag}, and the leverage of causal graphs to improve logical continuity~\citep{causalrag}. Additionally, retrieval strategies on graph KBs increasingly rely on heuristic path exploration, such as the topology-enhanced lightweight search~\citep{minirag}, the pruning via relational path retrieval~\citep{pathrag}, the utilization of personalized memory-inspired reasoning~\citep{hipporag, hipporag2}, and the adoption of beam search over proposition paths~\citep{proprag}.
Despite these advances, current GraphRAG approaches remain constrained by shallow retrieval, limiting their ability to perform deep searching over graph KBs.


\subsection{Agentic Retrieval-Augmented Generation}

RAG improves factual grounding by retrieving external knowledge~\citep{rag}, but single-round interaction is insufficient for complex reasoning tasks. Early advances focus on atomic-level improvements of RAG in query decomposition~\citep{decomposition}, query rewriting~\citep{queryrewrite,rqrag}, retrieval compression~\citep{recomp}, and selective retrieval decisions~\citep{slimplm}, which refine the retrieval process at a fine granularity. Beyond these, modular RAG systems~\citep{modularrag,flashrag,composerag} have been proposed to flexibly reconfigure retrieval and reasoning modules into composable pipelines. More recently, agentic approaches emerged, enabling LLMs to iteratively plan, retrieve, and reflect. Representative methods include reasoning–acting synergy in ReAct~\citep{react}, self-reflective retrieval in Self-RAG~\citep{selfrag}, test-time planning in PlanRAG~\citep{planrag}, and reinforcement-learned search agents in Search-o1~\citep{searcho1} and Search-r1~\citep{searchr1}. Subsequently, pioneering works~\citep{tog, tog2, gear, hybgrag} integrated structural graph knowledge for retrieval into the agentic RAG workflow to support the multi-step reasoning.

